\chapter{Voting Schemes}
\label{chap:votingschemes}
Voting schemes are used to conduct elections. 
Here, several voters submit ballots which are then aggregated by a voting authority. 
After that, the voting authority publishes the final tally of the election.

\todo{Include diagram}

We will now formalize this approach and iteratively extend it to fullfill several security considerations.

To motivate the different steps, we introduce a simple example: Electing a new president of the Computer Science Club at a university.
The club has the five members Alice (A), Bob (B), Charlie (C), Derpy (D) and Eve (E). 
Three of these members are running for the position, namely Alice, Bob and Eve.
In the election, each of the five clubs' members can vote for one candidate. 
It is possible to abstain from the vote.
% The votes will be counted by the voting authority ($VA$) Fred (F) who does not participate in the election.
In the end, the candidate that receives the most votes overall wins the election and becomes the next president.

In a first attempt at implementing this election, each of the members will write their chosen candidate on a piece of paper. 
The ballots will then be collected by Fred who counts how often each candidate was choosen and posts the result to a public bulletin board ($PBB$).
Fred does not participate in the election.

\todo{Naive diagram (Without standardized ballot format and with abstention by not casting a vote at all)}

One obvious problem with this approach is the lack of consistency of the ballots submitted by the members.
For example, it is not clear what to do if a voter misspells the name of a candidate on their ballot. 
We can combat this problem by standardizing the ballot format.

\section{Ballot Format}
In every voting scheme, the ballots submitted by the voters need to be chosen from specific choice space $C$. 
In our example, every voter now has to submit a ballot from $C:=\{(1,0,0), (0,1,0), (0,0,1), (0,0,0)\}$. 
Here a ballot $(a, b, e)$ signifies that the member gave $a$ votes to Alice, $b$ votes to Bob and $e$ Votes to Eve.

\todo{Example Voting Scheme with standardized Ballot format.}

Another benefit of standardized ballots is that the votes can now be counted automatically.
The new voting authority $VA$ (\todo{Robo-Fred} would be a little to infantile, right?) can aggregate the ballots by adding them component-wise via a function $agg$.
Then, they can publish the result of this computation.
The winner of the election can be determined by anyone based on this aggregated ballot.

In the current version of the voting scheme, votes are transmitted to the voting authority in plain.
This means that any party eavedropping this connection can see who the individual members voted for.
To prevent this, we need to add encryption.

\section{Confidentiality of Ballots}
To ensure that ballots are secret and nobody but the voters themself know what they voted for, ballots are encrypted by an encryption function $enc$ before being submitted.
We use an asymmetric encryption scheme for the encryption and decryption process as this allows users to encrypt ballots using a public key $k$ and without having to exchange a private key $\hat{k}$ beforehand.
Moreover, it is important to choose an additively homomorphic encryption scheme. 
With additive homomorphism, the voting authority can still aggregate the ballots by performing some computation $agg'$ on the encrypted ballots.
The result of the election can then be obtained by decrypting the output of $agg'$ using the corresponding decryption function $dec$ of the encryption scheme.
Since the encrypted ballots don't reveal any information about the plain ballots, it doesn't matter who obtains them. 
Therefore, the voters can submit their ballots to the public bulletin board and the voting authority can take them from there.

\todo{Ex. Voting Scheme with enc, dec}

A popular choice for an additively homomorphic encrytion scheme is Exponential ElGamal encryption. 
This will be covered in detail later on.

%If the voting authority is honest, confidentiality of the individual ballots is ensured in this new scheme.

Unfortunately, confidentiality is still not ensured in this new scheme.
The voting authority needs to have access to the private key of the used encrytion scheme to be able to decrypt the final tally of the election.
However, this means that a misbehaving voting authority can also decrypt individual votes. 
One approach to combat this problem is Threshhold encryption \todo{cite cryptographicVoting section 4.1}.
% Here, the public and private key of the encryption scheme are split into $n$ shares and distributed accross multiple authorities. 
% In our example, the members of the club could also be the authorities.
% Users can obtain the public key needed for encryption by combining the publicly available shares of this key using a combination function. %$comb$.
% However, decryption is only possible if at least $t$ of the authorities participate in the decryption process with their share of the private key $\hat{k}$.
% Additionally, confidentiality of individual ballots is only broken if at least $t$ authorities behave malicously and work together.

Exponential ElGamal is an encryption scheme that allows Threshold encryption.
Therefore it is still a good choice in this new setting.
% We will show how this is possible when introducing Exponential ElGamal.

Since the voting authority is not the focus of this thesis, we will asumme that it is honest. 
In that case there is no need for threshold encryption.
Note that an honest voting authority also eliminates the need for ensuring correct decryption of the tally.
If the voting authority was dishonest, they might just publish an invented tally that is favorable to them.
We can combat this problem using ZKPs \todo{cite cryptographicVoting section 4.3}.

% \todo{Ex. Voting Scheme with MPC encryption and decryption}

% Despite confidentiality of the ballots being ensured there remains another problem. 
% Currently, there is no guarantee of the integrity of the computed tally.
% A misbehaving authority might just publish an invented tally that is favorable to them.
% In the current version of the election scheme, we have no way of ensuring that the published tally is actually correct.
% This is what we change next.

% \section{Ensuring correct Decryption of the Tally}
% To ensure that the authorities decrypt the tally of the election correctly, we employ Zero Knowledge Proofs (ZKPs) \todo{Make acronym}. 
% The goal of ZKPs is to convince some party $X$ of the truth of a statement without revealing any confidential information to them.
% In our example this would mean convincing the general public that the final tally the authorities decrypted actually belongs to the aggregated encrypted tally.
% 
% For now we will look at ZKPs as Black-Box-Protocols $ZKP=(P,V)$. 
% Here, the proving algorithm $P$ takes some public inputs $I$ (also called instance) and some private inputs $W$ (also called witness) and produces a ZKP $\pi$.
% $\pi$ proves that the party who executed $P$ knows a corresponding witness $W$ for the instance $I$, even if only the instance is published.
% Anyone can verify the proof $\pi$ by running the verifying algorithm $V$ on $\pi$ and the instance $I$.

Even with an honest voting authority two problems that are directly connected remain.
In the current version of the election scheme, voters can submit ballots that are not in the choice space to gain an unfair advantage. 
In our example Eve could submit a vote $(0,0,100)$ to win the election. 

\todo{Ex. with Eve submitting $(0,0,100)$.}

Therefore, we need mechanisms to ensure the following:
\begin{enumerate}
  \item The encrypted ballot submitted by the voter actually belongs to the plain ballot they chose.
  \item The plain ballot chosen by the voter is actually part of the choice space.
\end{enumerate}

\section{Ensuring correct Ballot Format and Correct Encryption of Ballots}
In order to make sure that voters can only choose ballots from the choice space and then also encrypt exactly these ballots, we employ ZKPs \todo{Make acronym}. 

The goal of ZKPs is to convince some party of the truth of a statement without revealing any confidential information to them.
In our example, this would mean convincing the general public that the ballot a voter chose is part of the choice space and belongs to the encrypted ballot they submitted.
 
For now we will look at ZKPs as Black-Box-Protocols $zkp=(pro,ver)$. 
Here, the proving algorithm $pro$ takes some public inputs $I$ (also called instance) and some private inputs $W$ (also called witness) and produces a ZKP $\pi$.
$\pi$ proves that the party who executed $P$ knows a corresponding witness $W$ for the instance $I$, even though only the instance is published.
Anyone can verify the proof $\pi$ by running the verifying algorithm $ver$ on $\pi$ and the instance $I$.

In an election scheme this would mean that each voter $V$ runs the proving algorithm $pro$ to obtain a ZKP $\pi_V$.
Here, the witness is the plain ballot $b_V$ they chose and the instance is the corresponding encrypted ballot $b'_V$.
The member then publishes the encrypted ballot $b'_V$ along with the ZKP $\pi_V$ to the PBB.

Before aggregating the ballots, the voting authority runs $ver$ for each pair $(\pi_V, b'_V)$. 
The voting authority only accepts ballots with a valid ZKP $\pi_V$.

\todo{Ex. with ZKPs for ballot validity and correctness of encryption}

With this, we have reached the election scheme we will use throughout the rest of this thesis.

\todo{What is stopping voters from submitting multiple ballots to the public bulletin board?}

\section{Defining election schemes}

We can now define what an election scheme is formally. 
Here, we use asymmetric encryption schemes as presented in \todo{cite Katz and lindell, Section 12}. 
The voting scheme is based on the approach presented in \todo{cite fifty shades of ballot privacy}.

\begin{definition}[Choice space]
    A choice space is a non-empty set $C$.
\end{definition}

\begin{definition}[Public Bulletin Board]
    Let $C$ be a choice space. 
    Let $\mathcal{S}=(\text{Gen}(1^\eta), enc, dec)$ be a public-key (asymmetric) encryption scheme with:
    For any public key $pk$ generated by $\text{Gen}(1^\eta)$ and the corresponding message space $\mathcal{M}_{pk}$: $C\subseteq\mathcal{M}_{pk}$.
    Let $(pk, sk)$ be a fixed key pair generated by $\text{Gen}(1^\eta)$.
    Let $zkp=(pro,ver)$ be the Zero-Knowledge protocol allowing a voter to prove ballot validity.
    Then, a public bulleting board is $PBB=(C, pk, votes, b_{total}, publish, obtain)$.
  \todo{Publish and Obtain functionality}
\end{definition}

\begin{definition}[Voter]
  Let $PBB$ be a public bulleting board.
  Then, a voter is $V=(PBB, id, vote)$ with the algorithm vote defined in \ref{alg:vote}.
  \begin{Algorithmus} %Use the environment only if you want to place the algorithm similar to graphics from TeX
    \caption{Voter}
    \label{alg:voter}
    \begin{algorithmic}
    \Procedure{$vote$}{$PBB=()$}
      \State $b_V \gets C$
      \State $b'_V \gets enc_{pk}(b_V)$
      \State $\pi_V \gets pro(b_V, b'_V, pk, C)$
      \State $pub(b'_V, \pi_v)$

      \State \Comment $(a',c'a) \in \mathsf{HR}$ denotes that $a'$ is the parent of $a$
      \If{$\mathsf{parentHandled}\,\land(\mathcal{L}_\mathit{in}(a)=\emptyset\,\lor\,\forall l \in \mathcal{L}_\mathit{in}(a): \mathsf{visited}(l))$}
      \State $\mathsf{visited}(a) \gets \text{true}$
      \State $\mathsf{writes}_\circ(a,v_e) \gets
      \begin{cases}
      \mathsf{joinLinks}(a,v_e) & \abs{\mathcal{L}_\mathit{in}(a)} > 0\\
      \mathsf{writes}_\circ(p,v_e)
      & \exists p: (p,c,a) \in \mathsf{HR}\\
      (\emptyset, \emptyset, \emptyset, false) & \text{otherwise}
      \end{cases}
      $
      \If{$a\in\mathcal{A}_\mathit{basic}$}
        \State \Call{HandleBasicActivity}{$a$,$v_e$}
      \ElsIf{$a\in\mathcal{A}_\mathit{flow}$}
        \State \Call{HandleFlow}{$a$,$v_e$}
      \ElsIf{$a = \mathsf{process}$} \Comment Directly handle the contained activity
        \State \Call{HandleActivity}{$a'$,$v_e$}, $(a,\bot,a') \in \mathsf{HR}$
        \State $\mathsf{writes}_\bullet(a) \gets \mathsf{writes}_\bullet(a')$
      \EndIf
      \ForAll{$l \in \mathcal{L}_\mathit{out}(a)$}
        \State \Call{HandleLink}{$l$,$v_e$}
      \EndFor
      \EndIf
    \EndProcedure
    \end{algorithmic}
  \end{Algorithmus}
\end{definition}

\begin{definition}[Honest Voting Authority]
  
\end{definition}

\begin{definition}[Election Scheme]
  Let $C$ be a non-empty choice space, let $\mathcal{S}=(C, \text{Gen}(1^\eta), enc:C\to\tilde{C}, dec:\tilde{C}\to C)$ be an encryption scheme (\todo{Explain/cite}) with $C\subseteq X$.
  
\end{definition}

\begin{definition}[Voting system]
  Let $C$ be a non-empty choice space, let $\mathcal{S}=(C, \text{Gen}(1^\eta), enc:C\to\tilde{C}, dec:\tilde{C}\to C)$ be an encryption scheme (\todo{Explain/cite}) with $C\subseteq X$.
  Let $agg:\tilde{C}^*\to\tilde{C}$ be an aggregation function. Then, the voting system $VS=(\mathcal{S}, agg)$
\end{definition}






\section{Voting schemes}

In an election, multiple voters can submit ballots and a voting authority computes the result of the election from these ballots. 
To ensure that nobody, including the voting authority, learns anything about the ballots chosen by the voters, the voters only submit the encrypted ballots.

In a voting scheme, a voter needs to execute the following steps:
\begin{itemize}
  \item Choose a ballot $b$ from an allowed range of ballots called choice space $C$.
  \item Encrypt the ballot using some (semi-)homomorphic encryption scheme.
  \item Publish their ballot to a public bulletin board.
\end{itemize}
The voting authority does the following computations:
\begin{itemize}
  \item Do some computations on the encrypted ballots $\tilde{b}_1,\tilde{b}_2,\dots,\tilde{b}_n$ submitted by the voters using an aggregation function $agg$.
  \item Decrypt the aggregated ballot $\tilde{b}_total$.
  \item Output the final aggregated tally $b_{total}$.
\end{itemize}

Unfortunately, this approach still has some security issues. Specifically, the following two questions are unanswered by the approach:
\begin{enumerate}
  \item How can we convince the voting authority that the submitted encrypted ballots $\tilde{b}_i$ belong to a valid ballots $b_i$ from the choice space $C$ without obtaining any information about the plain ballots $b_i$?
  \item How can we ensure that the voting authority only decrypts the final tally $\tilde{b}_{total}$ and not any of the individual votes $\tilde{b}_i$?
\end{enumerate}

This thesis will focus on the first question. 
We can solve the second question using distributed decryption where multiple trustees choose their own public keys and decryption is only possible by combining all of them. 
Adida et. al. present this solution for Exponential ElGamal encryption %\cite{helios}.

\section{Proving Ballot Validity}
A voter $V_i$ can employ a Zero-Knowledge-Proof (ZKP) to convince the voting authority in a voting system that the encrypted ballot $\tilde{b}_i$ they submitted belongs to a plain ballot $b_i$ from the choice space $C$. 
By using a ZKP, the voter does not disclose any information about the plain ballot $b_i$ to the voting authority.

We can formalize this as follows:
\begin{definition}[Voter]
  Let $C$ be a non-empty choice space, let $\mathcal{S}=(C, \text{Gen}(1^\eta), enc:C\to\tilde{C}, dec:\tilde{C}\to C)$ be an encryption scheme (\todo{Explain/cite}) with $C\subseteq X$.
  Then, a voter $V$ is defined by the following algorithm:
  
\end{definition}

\begin{definition}[Voting system]
  Let $C$ be a non-empty choice space, let $\mathcal{S}=(C, \text{Gen}(1^\eta), enc:C\to\tilde{C}, dec:\tilde{C}\to C)$ be an encryption scheme (\todo{Explain/cite}) with $C\subseteq X$.
  Let $agg:\tilde{C}^*\to\tilde{C}$ be an aggregation function. Then, the voting system $VS=(\mathcal{S}, agg)$
\end{definition}

For now, 

On closer inspection we can split this into two seperate problems:
\begin{enumerate}
  \item How can we convince the voting authority that the voter $i$ knows a plain ballot $b_i$ that belongs to the encrypted one $\tilde{b}_i$?
  \item How can we convince the voting authority that this plain ballot $b_i$ is valid?
\end{enumerate}

To convince the voting authority of the truth of these two statements, we employ Zero-Knowledge-Proofs.